\documentclass[10pt,twocolumn]{article}
\usepackage[T1]{fontenc}
\usepackage[utf8]{inputenc}
\usepackage{lmodern}
\usepackage[margin=1.8cm,top=2.4cm,bottom=2cm,columnsep=0.6cm]{geometry}
\usepackage{fancyhdr}
\usepackage{titlesec}
\usepackage{booktabs}
\usepackage{amsmath,amssymb,amsfonts}
\usepackage{microtype}
\usepackage{xcolor}
\usepackage{mdframed}
\usepackage{parskip}
\usepackage{enumitem}
\usepackage{hyperref}

\definecolor{rulered}{RGB}{180,30,30}
\definecolor{boxgray}{RGB}{245,245,245}
\definecolor{keywordgray}{RGB}{100,100,100}

\pagestyle{fancy}
\fancyhf{}
\fancyhead[L]{\textsc{\small Claude Mag}}
\fancyhead[C]{\small\textit{Machine Learning Research Digest}}
\fancyhead[R]{\small 2026-08-17}
\fancyfoot[C]{\small\thepage}
\renewcommand{\headrulewidth}{0.8pt}

\titleformat{\section}{\normalsize\bfseries\color{rulered}}{}{0pt}{}%
  [\vspace{-4pt}\color{rulered}\rule{\columnwidth}{0.4pt}]
\titlespacing{\section}{0pt}{8pt}{4pt}

\hypersetup{colorlinks=true,urlcolor=rulered,linkcolor=black}

\begin{document}
\setlength{\parindent}{0pt}
\setlength{\parskip}{4pt}

{\small\color{keywordgray}\textbf{Keywords:}
  diffusion language models \textbullet{}
  masked decoding \textbullet{}
  reasoning order \textbullet{}
  test-time inference}\par\vspace{4pt}

{\LARGE\bfseries\raggedright
  Answer First, Reason Later:
  Commitment Order in Diffusion LLMs}\par\vspace{4pt}

{\large\itshape\raggedright
  Unconstrained Masked Decoding Commits Final Answers Before Reasoning Completes,
  Collapsing to Answer-Only Outputs on Up to 90\% of Problems}\par\vspace{6pt}

{\small\textbf{By} Jewon Yeom, Jaewon Sok, Seonghyeon Park, Jeongjae Park,
  Hwiyeong Lee, Taesup Kim}\par\vspace{2pt}

{\small\color{keywordgray}arXiv:2608.05687 \textbullet{} Preprint, August 2026}
\par\vspace{2pt}\noindent{\color{rulered}\rule{\columnwidth}{0.6pt}}\vspace{6pt}

Masked diffusion language models promise order-free generation as a core
advantage over autoregressive decoding. A new study from Seoul National University
shatters that assumption for reasoning tasks: logging every token commitment in
LLaDA-8B on GSM8K reveals that unconstrained decoding commits the final answer
at only 15--24\% of the trajectory, before half the reasoning region is written,
collapsing to answer-only outputs on up to 90\% of problems as canvas size grows.

\begin{mdframed}[backgroundcolor=boxgray,linecolor=rulered,linewidth=1pt,
    innerleftmargin=6pt,innerrightmargin=6pt,
    innertopmargin=6pt,innerbottommargin=6pt]
\textbf{\small AT A GLANCE}\par\vspace{4pt}
\begin{description}[leftmargin=0pt,labelsep=4pt,itemsep=2pt,parsep=0pt,
    font=\small\bfseries]
  \item[Problem:] Diffusion LLMs generate tokens in arbitrary order, but reasoning tasks require
    intermediate steps to precede the final answer — an ordering constraint the model violates.
  \item[Why it matters:] dLLMs are increasingly proposed as reasoning engines; if their decoding
    procedure systematically skips reasoning, the architecture is unreliable for this use case.
  \item[Method:] Log every token commitment (masked $\to$ unmasked transition) during LLaDA-8B
    decoding on GSM8K; analyze commitment position as a fraction of the decoding trajectory.
  \item[Result:] Pure decoding commits the answer at 15--24\% of the trajectory; constrained
    ordering recovers substantial accuracy without retraining.
  \item[Evidence:] Collapse to answer-only outputs reaches 90\% of problems at large canvas size;
    commitment trajectory visualizations confirm answer-first ordering.
\end{description}
\end{mdframed}

\section{The Big Picture}

Masked diffusion language models (dLLMs) generate text by starting from a fully
masked token sequence and iteratively unmasking tokens over multiple diffusion steps.
Unlike autoregressive (AR) models, which are constrained to generate left-to-right,
dLLMs may unmask any position in any order. Proponents argue this flexibility enables
superior global coherence and parallel generation.

For reasoning tasks, however, this flexibility has a fatal cost. A reasoning trace
has inherent logical structure: the intermediate steps should be generated \emph{before}
the final answer, because the answer should be justified by preceding computation, not
the reverse. When a dLLM commits the answer token before writing the reasoning steps,
it is, in effect, guessing — and the reasoning region becomes post-hoc decoration for
a decision already made.

\section{Why Existing Approaches Fall Short}

Standard dLLM decoding (``pure'' diffusion) selects which tokens to unmask at each
step based on model confidence, without any constraint on the semantic identity of
the tokens being committed. The model may unmask the final answer position early
because the answer is high-confidence given the problem statement alone — before
the chain-of-thought region provides any additional information.

Prior work on improving dLLM reasoning has focused on training objectives (RLVR,
supervised fine-tuning on CoT data) or on post-hoc verification. None has examined
the \emph{order} in which tokens are committed during decoding, which this paper
identifies as the proximate cause of failure.

\section{Core Mechanism}

The authors introduce \textbf{commitment order logging}: at each diffusion step, they
record which output position transitions from masked to unmasked, forming a complete
commitment trajectory. For a generation with $T$ positions and $S$ decoding steps:

\begin{itemize}[itemsep=2pt,parsep=0pt]
  \item Each step $s$ produces a set $C_s \subseteq [T]$ of newly unmasked positions.
  \item The answer region $\mathcal{A} \subset [T]$ and reasoning region $\mathcal{R} \subset [T]$
    are identified by position in the output template.
  \item The ``answer commitment step'' $s^* = \min\{s : C_s \cap \mathcal{A} \neq \emptyset\}$
    is the fraction $s^*/S$ of the trajectory at which the final answer is first committed.
\end{itemize}

The \textbf{constrained decoding} fix enforces that positions in $\mathcal{A}$ remain
masked until a threshold fraction $\tau$ of the trajectory has elapsed, releasing the
answer region only after substantial reasoning has been committed.

\section{Technical Formulation}

Let $\mathbf{x}_0$ be the target output and $\mathbf{x}_s$ the partially masked sequence
at step $s$, with $\mathbf{x}_S$ fully masked. Pure reverse diffusion selects tokens
to unmask according to:
\[
p_\theta(\mathbf{x}_{s-1} \mid \mathbf{x}_s) \propto \prod_{i \in [T]} p_\theta(x_i \mid \mathbf{x}_s)^{\mathbb{1}[x_s^{(i)} = \texttt{[MASK]}]}
\]
Commitment at step $s$ selects the top-$k$ unmasked positions by $p_\theta(x_i \mid \mathbf{x}_s)$.

Constrained decoding modifies the selection by adding a region-based gate:
\[
p^{\mathrm{con}}_\theta(x_i \mid \mathbf{x}_s) = p_\theta(x_i \mid \mathbf{x}_s) \cdot \mathbb{1}\!\left[i \notin \mathcal{A} \;\text{ or }\; \frac{s}{S} > \tau\right]
\]
where $\tau \in (0,1)$ is the threshold fraction. No change to model weights is required;
the constraint is purely at inference time.

\section{Learning or Inference Procedure}

No training is involved. Both pure decoding and constrained decoding are evaluated
at inference time on LLaDA-8B, a masked diffusion language model with 8 billion
parameters. GSM8K provides 1,319 test-set math word problems as the evaluation
benchmark. The commitment threshold $\tau$ is tuned on a small held-out validation set.

\section{What the Guarantee Says}

No formal guarantee is provided. The empirical claim is that answer-first commitment
is the dominant failure mode of pure dLLM decoding on reasoning tasks, and that
constraining the answer region during early decoding consistently recovers accuracy.
The correlation between canvas size and collapse rate (peaking at 90\% answer-only
outputs) establishes that the failure is deterministic as generation length grows,
not stochastic.

\section{Experimental Findings}

\begin{itemize}[itemsep=2pt,parsep=0pt]
  \item \textbf{Commitment timing:} Pure decoding commits the answer at 15--24\% of
    the decoding trajectory on average.
  \item \textbf{Collapse rate:} Up to 90\% of problems produce answer-only outputs
    (zero reasoning steps) as canvas size grows.
  \item \textbf{Constrained decoding:} Enforcing reasoning-before-answer commitment
    recovers substantial GSM8K accuracy with threshold $\tau = 0.5$.
  \item \textbf{Scale sensitivity:} The gap between constrained and unconstrained
    decoding widens with output length, confirming scale-dependent failure.
\end{itemize}

\section{Ablations and Interpretation}

\textbf{Why does pure decoding answer first?}
The answer token has high local confidence given the problem statement alone
(the model ``knows'' plausible answer formats). The reasoning region is globally
coherent only \emph{after} intermediate steps are committed — but those steps are
low-confidence before being written, so the decoding procedure deprioritizes them.

\textbf{Why does collapse worsen with canvas size?}
Longer outputs have more reasoning positions that must be committed before the
answer is released. As the reasoning region grows, the denominator of the
commitment problem grows — making early, spurious answer commitment
increasingly likely under unconstrained selection.

\textbf{Threshold sensitivity:} Accuracy peaks around $\tau = 0.4$--$0.6$ with
a shallow optimum; extreme values ($\tau \to 0$ or $\tau \to 1$) recover the
failure modes of pure decoding or fully blocked answer generation respectively.

\section*{Reference}

Jewon Yeom, Jaewon Sok, Seonghyeon Park, Jeongjae Park, Hwiyeong Lee, Taesup Kim.
\textbf{Answer First, Reason Later: Commitment Order in Diffusion LLMs}.
arXiv:2608.05687, August 2026.
\url{https://arxiv.org/abs/2608.05687}

\end{document}
